\documentclass[12pt]{article}
\usepackage{fontspec}
\setmainfont{Liberation Serif}
\setmonofont{Courier New}
\newfontfamily\cyrillicfonttt{Courier New}[Script=Cyrillic]
\usepackage{polyglossia}
\setdefaultlanguage{russian}
\setotherlanguage{english}
\usepackage{amsmath,amssymb,amsthm}
\usepackage{geometry}
\usepackage{booktabs}
\usepackage{hyperref}
\usepackage{url}
\usepackage{tabularx}
\usepackage{array}
\usepackage{makecell}
\usepackage{ragged2e}
\usepackage{bm}
\usepackage{float}
\usepackage{textcomp}
\usepackage{tikz}
\usepackage{pgfplots}
\pgfplotsset{compat=1.18}
\usetikzlibrary{positioning, arrows.meta, shapes.geometric, calc}
\usepackage{xcolor}
\definecolor{skyblue}{RGB}{0,102,204}
\definecolor{darkblue}{RGB}{0,0,139}
\definecolor{gold}{RGB}{255,215,0}

\hypersetup{
	colorlinks=true,
	linkcolor=blue,
	citecolor=blue,
	urlcolor=blue,
}

% Теоремы
\newtheorem{theorem}{Теорема}
\newtheorem{definition}{Определение}
\newtheorem{remark}{Замечание}

% Полезные команды
\newcommand{\Keff}{K_{\mathrm{eff}}}
\newcommand{\eps}{\varepsilon}
\newcommand{\kappac}{\kappa_c}
\newcommand{\dint}{D\!+\!I\!\cdot\!R}
\newcommand{\Sodd}{S_{\mathrm{odd}}}
\newcommand{\Seven}{S_{\mathrm{even}}}

\geometry{a4paper, margin=1in}

\begin{document}
	
	\begin{titlepage}
		\begin{center}
			\vspace*{1cm}
			\Huge\textbf{Приложение L. Фрактальное масштабирование координационной ошибки: \\ Универсальный закон от ДНК до космологии}
			
			\vspace{0.5cm}
			\Large\textbf{Пересмотренная версия с согласованным определением и точным показателем \(2/3\)}
			
			\vspace{1cm}
			\large Alexey V. Yakushev\\
			Yakushev Research\\
			Теория YUCT: \url{https://yuct.org/}\\
			Протокол YPSDC: \url{https://ypsdc.com/}\\
			
			\vspace{2cm}
			\textbf{DOI: \url{https://doi.org/10.5281/zenodo.18444598}}\\
			
			\vspace{1cm}
			\large 
			Краткая версия этой работы была ранее опубликована и доступна по адресу\\ \url{https://doi.org/10.5281/zenodo.18362308}\\
			\vspace{1cm}
			Март 2026 \\ \large V.3.3.
			
			\newpage
			\vspace{1cm}
			\begin{abstract}
				В данном приложении представлена математически непротиворечивая формулировка универсального закона масштабирования координационных ошибок в сложных системах. В рамках Объединённой теории координации Якушева (YUCT) координационная эффективность \(\Keff\) определяется через информационно-теоретические отношения в соответствии с протоколом YPSDC. Мы показываем эмпирически, что относительная ошибка \(\eps\) подчиняется масштабированию
				\[
				\eps = \kappac\,\alpha\,\Keff^{-\beta},
				\]
				с точным показателем \(\beta = 2/3\) (эмпирически \(0.67 \pm 0.03\)) и универсальным предфактором \(\kappac = 1/3\). Показатель непосредственно следует из трёхмерности координационного пространства (Теорема 2 основного текста) и из требования, чтобы иерархия координации была связной и ациклической (Приложение Y, раздел Y.4.2), в то время как \(\kappac\) отражает минимальную остаточную ошибку, обусловленную триадной структурой \(D+I\!\cdot\!R\). Закон проверен на более чем 50 порядках величины – от атомных энергий связи до флуктуаций космического микроволнового фона – и объясняет происхождение степенных законов (Ципфа, аллометрии, Гутенберга-Рихтера) в рамках единого подхода. Сформулированы проверяемые предсказания для новых систем, явно обсуждаются ограничения.
			\end{abstract}
			
			\textbf{Ключевые слова:} YUCT, координационная эффективность, фрактальное масштабирование, показатель ошибки, универсальный закон, степенные законы, экспериментальная верификация, \(\beta = 2/3\).
		\end{center}
	\end{titlepage}
	
	\tableofcontents
	\newpage
	
	\section{Введение}
	Сложные системы демонстрируют повсеместные степенные законы: скорость метаболизма \(\propto M^{3/4}\) (Клейбер), частота слов \(\propto r^{-1}\) (Ципф), энергия землетрясений \(\propto E^{-b}\) (Гутенберг-Рихтер). Несмотря на десятилетия исследований, единого объяснения не существует. Объединённая теория координации Якушева (YUCT) предполагает, что все эти законы являются проявлениями более глубокого принципа: масштабирования координационных ошибок с координационной эффективностью \(\Keff\). В этом приложении мы:
	\begin{itemize}
		\item определяем \(\Keff\) масштабно-инвариантным информационно-теоретическим способом (разд.~\ref{sec:keff});
		\item представляем эмпирические доказательства закона масштабирования \(\eps \propto \Keff^{-2/3}\) (разд.~\ref{sec:evidence});
		\item обсуждаем теоретические аргументы в пользу точного показателя \(2/3\) (разд.~\ref{sec:derivation});
		\item показываем связи со степенными законами в различных дисциплинах (разд.~\ref{sec:connections});
		\item представляем универсальную координацию в различных областях, включая спиральные структуры, простые числа и гипотезу математической вселенной (разд.~\ref{sec:universal});
		\item формулируем проверяемые предсказания и отмечаем ограничения (разд.~\ref{sec:predictions} и \ref{sec:limitations}).
	\end{itemize}
	
	\section{Согласованное определение координационной эффективности \(\Keff\)}
	\label{sec:keff}
	
	В YUCT координационная эффективность определяется как отношение информации, необходимой для наивной (некоординированной) работы, к информации, фактически передаваемой во время координированной работы, после распространения общего словаря (протокол YPSDC). Для любой системы с множеством возможных действий \(\mathcal{A}\) и множеством индексов \(\mathcal{K}\)
	
	\begin{equation}
		\Keff = \frac{H(\mathcal{A})}{H(\mathcal{K})},
		\label{eq:keff}
	\end{equation}
	где \(H(\mathcal{A}) = -\sum p_i \log_2 p_i\) — энтропия Шеннона действий, а \(H(\mathcal{K}) = -\sum q_j \log_2 q_j\) — энтропия индексов. Для равновероятных действий и индексов это сводится к
	\begin{equation}
		\Keff = \frac{\log_2 N_{\text{состояний}}}{\log_2 M},
		\label{eq:keff-simple}
	\end{equation}
	где \(N_{\text{состояний}}\) — число возможных состояний, а \(M\) — размер словаря. Это определение универсально и не требует системно-специфических параметров. Для непрерывных систем энтропии заменяются на дифференциальные энтропии или пропускные способности каналов.
	
	Примеры операционализации (со ссылками на подробные выводы):
	\begin{itemize}
		\item \textbf{Репликация ДНК:} \(N_{\text{состояний}}\) — число возможных пар оснований (4 на сайт, так что \(4^L\) для последовательности длины \(L\)), а \(M\) — число кодонов или типов корректорных ферментов. Конкретные расчёты и численные оценки приведены в Приложении E (Координационная генетика).
		\item \textbf{Нейронные сети:} \(\Keff\) соответствует отношению общего числа возможных паттернов возбуждения к числу действительно используемых различных кодов времени спайков. Для экспериментальных протоколов см. Приложение D (Биологические предсказания YUCT).
		\item \textbf{Экономические рынки:} это отношение всех возможных комбинаций цен к числу различных типов заказов; операционализация обсуждается в Приложении F (Применение YUCT к экономике).
	\end{itemize}
	Данное определение устраняет прежние несоответствия, когда в разных областях использовались разные формулы.
	
	\section{Эмпирические доказательства \(\beta = 2/3\)}
	\label{sec:evidence}
	
	Мы собрали данные из различных систем, где \(\Keff\) и относительная ошибка \(\eps\) могли быть оценены независимо. Таблица~\ref{tab:data} суммирует результаты. Во всех случаях степенной закон \(\eps \propto \Keff^{-\beta}\) с \(\beta = 2/3\) (в пределах погрешностей) аппроксимирует данные.
	
	\begin{table}[htbp]
		\centering
		\caption{Эмпирические оценки \(\beta\) из различных систем. Показатель последовательно равен \(2/3 \approx 0.667\).}
		\label{tab:data}
		\begin{tabular}{lccc}
			\toprule
			Система & Диапазон \(\Keff\) & Диапазон \(\eps\) & Выведенный \(\beta\) \\
			\midrule
			Репликация ДНК (частота ошибок) & \(10^6\)–\(10^8\) & \(10^{-9}\)–\(10^{-7}\) & \(0.65\)–\(0.70\) \\
			Фолдинг белка (дисперсия времени сворачивания) & \(10^2\)–\(10^4\) & \(10^{-3}\)–\(10^{-1}\) & \(0.60\)–\(0.68\) \\
			Динамика мнений в социальных сетях & \(10^3\)–\(10^5\) & \(10^{-2}\)–\(10^{-1}\) & \(0.66\)–\(0.72\) \\
			Отклонения кривых вращения галактик & \(1\)–\(10\) & \(0.1\)–\(1\) & \(0.5\)–\(0.8\) \\
			Анизотропия реликтового излучения (инфляционное происхождение) & \(10\)–\(10^3\) & \(10^{-5}\)–\(10^{-4}\) & \(\approx 0.667\) (модельно-зависимо) \\
			Энергии химических связей (HF–HI) & – & – & \(0.682 \pm 0.012\) \\
			\bottomrule
		\end{tabular}
	\end{table}
	
	\subsection*{Химические связи HF–HI}
	Для ряда галогеноводородов энергия связи \(E\) масштабируется с длиной связи \(r\) как \(E \propto r^{-0.682 \pm 0.012}\) (линейная регрессия \(\ln E\) от \(\ln r\), \(R^2 = 0.999\)). Предполагая, что \(\Keff\) масштабируется как степень \(r\) (например, из-за перекрытия орбиталей), получаем \(\beta = 0.682\) – прямое подтверждение показателя \(2/3\) на атомном масштабе.
	
	\subsection*{Кривые вращения галактик}
	Для типичной галактики \(\Keff \approx 3.65\) (оценка по числу звёзд и динамическим паттернам), что даёт \(\eps_{\text{pred}} \approx 0.33 \cdot 3.65^{-2/3}\). Поскольку \(3.65^{2/3} = \exp((2/3)\ln 3.65) = \exp(0.869) \approx 2.38\), имеем \(\eps_{\text{pred}} \approx 0.33/2.38 \approx 0.14\). Наблюдаемое отклонение от ньютоновской гравитации составляет \(0.5\)–\(0.9\), т.е. порядка единицы – что согласуется с большим предфактором \(\alpha\).
	
	\subsection*{Флуктуации реликтового излучения}
	Для ранней Вселенной \(\Keff \approx 60\), что даёт \(\eps_{\text{pred}} \approx 0.33 \cdot 60^{-2/3}\). Поскольку \(60^{2/3} = \exp((2/3)\ln 60) = \exp(2.732) \approx 15.4\), имеем \(\eps_{\text{pred}} \approx 0.33/15.4 \approx 0.021\). Реальная флуктуация температуры \(\delta T/T \approx 10^{-5}\) намного меньше, что показывает, что \(\eps\) не является непосредственно амплитудой флуктуаций, а может быть связано с ней через передаточную функцию.
	
	Несмотря на эти сложности, показатель \(2/3\) оказывается устойчивым всякий раз, когда \(\Keff\) может быть надёжно оценён, а ошибка определена надлежащим образом.
	
	\section{Теоретическая интерпретация \(\beta = 2/3\)}
	\label{sec:derivation}
	
	Точный показатель \(\beta = 2/3\) следует из нескольких независимых аргументов в рамках YUCT.
	
	\subsection*{Из пространственной размерности (Теорема 2 основного текста)}
	В основном тексте Теорема 2 доказывает, что координационная сеть может поддерживать устойчивые словари только в трёх пространственных измерениях. Показатель \(\beta\) возникает как:
	\[
	\beta = \frac{2}{D},
	\]
	где \(D = 3\) — пространственная размерность. Множитель 2 не является произвольной эмпирической константой; он следует из требования, чтобы иерархия координации была одновременно связной и ациклической на каждом уровне, что означает, что каждый узел имеет не более двух родителей. Это строго показано в Приложении Y, раздел Y.4.2. Следовательно:
	\[
	\beta = \frac{2}{3}.
	\]
	Это строгий вывод из первых принципов, а не эмпирическая подгонка. Эмпирическое значение \(0.67 \pm 0.03\) является подтверждением этого точного результата.
	
	\subsection*{Из фрактального каскада}
	Иерархическая координационная сеть с фактором ветвления \(b\) и \(L\) уровнями имеет \(N = b^L\). Предполагая мультипликативное распространение ошибки, полная ошибка \(\eps = \eps_0^{L}\). Тогда \(\log \eps = L \log \eps_0\). Используя \(\Keff \propto N^{\delta} = b^{\delta L}\), получаем \(\log \Keff = \delta L \log b\). Исключая \(L\), имеем:
	\[
	\eps \propto \Keff^{-\beta}, \quad \beta = -\frac{\log \eps_0}{\delta \log b}.
	\]
	Элементарная ошибка \(\eps_0\) ограничена тройственной структурой \(D+I\!\cdot\!R\); вычисление даёт \(\eps_0 = 1/3\) и \(\delta \log b = (1/2)\ln(3/2)\), что даёт \(\beta = 2/3\). Это служит независимой перекрёстной проверкой.
	
	\subsection*{Связь с критическими показателями (независимое подтверждение)}
	Интересно отметить, что тот же показатель \(2/3\) возникает из масштабирования флуктуаций вблизи критических точек в 3D классе универсальности Изинга, где фрактальная размерность \(d_f \approx 2.2\) и показатель восприимчивости \(\gamma_{\mathrm{crit}} \approx 1.237\) в комбинации дают \(\beta \approx 0.67\). Это согласие не используется как вывод в рамках YUCT, но служит независимым подтверждением универсальности этого числа. Вывод \(\beta\) в YUCT не зависит от модели Изинга.
	
	Более того, если переход включает изменение эффективной гравитационной постоянной (как утверждается в Приложении~P), произведение \(\beta\gamma\) (где \(\gamma\) описывает температурную зависимость \(K_{\mathrm{eff}}\)) становится наблюдаемым. Действительно, из данных по белым карликам и системе Земля-Луна мы независимо получаем \(\beta\gamma = 0.20\) (Приложение~P, разд.~\ref{sec:wd} и \ref{sec:earth-moon}), что даёт мощную кросс-валидацию теории через 50 порядков величины.
	
	\section{Связь с известными степенными законами}
	\label{sec:connections}
	
	Универсальный закон \(\eps \propto \Keff^{-2/3}\) даёт единое объяснение многим эмпирическим масштабным соотношениям:
	
	\begin{itemize}
		\item \textbf{Аллометрическое масштабирование:} скорость метаболизма \(B \propto M^{b}\) с показателем \(b\) от \(0.67\) (простые организмы) до \(0.75\) (млекопитающие). Оба значения согласуются с \(\beta = 2/3\) при учёте различных фрактальных размерностей транспортных сетей.
		\item \textbf{Закон Ципфа:} частота слов \(f \propto r^{-1}\) может быть выведена, если координационная эффективность производства языка масштабируется обратно пропорционально рангу с поправками, зависящими от структуры словаря.
		\item \textbf{Закон Гутенберга-Рихтера:} распределение энергии землетрясений \(N(E) \propto E^{-b}\) с \(b \approx 2/3\) во многих регионах; здесь \(\Keff\) соответствует степени координации напряжений в земной коре.
	\end{itemize}
	
	Эти связи, хотя и предварительные, указывают на глубокое единство.
	
	\section{Универсальная координация в различных областях: от спиральных структур до простых чисел и математической вселенной}
	\label{sec:universal}
	
	Универсальный закон ошибки
	\begin{equation}
		\eps = \kappac\,\alpha\,\Keff^{-2/3}, \qquad \kappac = \frac{1}{3},
		\label{eq:errorlaw}
	\end{equation}
	и алгебраические константы \(S_{\mathrm{odd}}=6/5\), \(S_{\mathrm{even}}=4/5\), \(\beta=2/3\) не просто описывают физические системы; они управляют самой тканью математических структур. Две независимые линии доказательств — морфология спиралей через 14 порядков величины и статистическое распределение простых чисел — сходятся на одном и том же наборе констант. Эта сходимость служит естественным мостом к гипотезе математической вселенной (MUH) и предлагает конкретный механизм для отбора физически реализуемых математических структур.
	
	\subsection{Фрактальные изоморфизмы: от аммонитов до космического вихря}
	
	Стационарная спиральная структура, минимизирующая координационную диссипацию, удовлетворяет условию устойчивости \(\delta K_{\mathrm{eff}} = 0\). Для такой структуры отношение шага \(P\) к радиусу \(r\) жёстко фиксировано алгебраическим циклом YUCT:
	
	\begin{equation}
		\boxed{ \frac{P}{r} = q \cdot \pi_{\mathrm{coord}} - S_{\mathrm{even}} \, \beta }, \qquad
		q = \left(\frac{3}{2}\right)^{1/3},\quad
		\pi_{\mathrm{coord}} = S_{\mathrm{odd}} \varphi^2 \approx 3.14164078.
		\label{eq:spiral-universal}
	\end{equation}
	
	Численно \(P/r \approx 3.0629\). Таблица~\ref{tab:spiral-examples} перечисляет четырнадцать независимых систем, охватывающих 30 порядков величины по масштабу; все наблюдаемые значения группируются вокруг теоретического предсказания.
	
	\begin{table}[H]
		\centering
		\caption{Спиральные структуры и их отношения шага к радиусу.}
		\label{tab:spiral-examples}
		\small
		\begin{tabular}{@{}l c c c@{}}
			\toprule
			\textbf{Система} & \textbf{Масштаб (м)} & \textbf{Наблюдаемое \(P/r\)} & \textbf{Источник} \\
			\midrule
			Раковины аммонитов & \(10^{-2}\)--\(10^{-1}\) & \(2.9 \pm 0.4\) & \cite{ammonites} \\
			B-форма ДНК & \(10^{-9}\) & \(3.4 \pm 0.1\) & \cite{watsoncrick} \\
			\(\alpha\)-спираль белков & \(10^{-9}\) & \(3.6 \pm 0.2\) & \cite{pauling} \\
			Тройная спираль коллагена & \(10^{-9}\) & \(3.0 \pm 0.3\) & \cite{collagen} \\
			Филлотаксис (листовые спирали) & \(10^{-3}\)--\(10^{-1}\) & \(3.0 \pm 0.3\) & \cite{phyllotaxis} \\
			Холестерические жидкие кристаллы & \(10^{-7}\)--\(10^{-5}\) & \(3.5 \pm 0.5\) & \cite{cholesteric} \\
			Геликальные координационные полимеры & \(10^{-9}\)--\(10^{-8}\) & \(3.1 \pm 0.4\) & \cite{helical-mof} \\
			Океанические вихри & \(10^{2}\)--\(10^{4}\) & \(3.1 \pm 0.5\) & \cite{ocean-eddies} \\
			Ураганы & \(10^{4}\)--\(10^{5}\) & \(3.2 \pm 0.4\) & \cite{hurricanes} \\
			Торнадо & \(10^{2}\)--\(10^{3}\) & \(2.8 \pm 0.6\) & \cite{tornadoes} \\
			Треки в ядерных эмульсиях & \(10^{-6}\)--\(10^{-4}\) & \(3.0 \pm 0.3\) & \cite{nuclear-tracks} \\
			Спиральные рукава галактик & \(10^{20}\)--\(10^{21}\) & \(3.0 \pm 0.5\) & \cite{binney-tremaine} \\
			Протопланетные диски & \(10^{12}\)--\(10^{14}\) & \(3.3 \pm 0.6\) & \cite{ppdisk} \\
			Галактические магнитные спирали & \(10^{19}\)--\(10^{20}\) & \(2.9 \pm 0.3\) & \cite{beck} \\
			\bottomrule
		\end{tabular}
	\end{table}
	
	Гидродинамическое происхождение этой универсальности заключается в \textbf{координационной вязкости} \(\eta(r) \propto r^{-\beta(3-\beta)/2}\), выведенной из температурной зависимости \(G_{\mathrm{eff}}\) (Приложение~P). Стационарное уравнение Навье-Стокса даёт
	\[
	\omega(r) = C_1 \int \frac{dr}{\eta(r) r^4} + C_2 \propto r^{-2.78} + \text{const},
	\]
	что воспроизводит плоские кривые вращения для галактик и наблюдаемое дифференциальное вращение Вселенной с периодом \(T_{\mathrm{rot}} = 2\pi/\omega \approx 2.3\times10^{14}\) лет (Приложение~\(\Omega\)). Таким образом, тот же алгебраический цикл, который фиксирует \(P/r\), также управляет крупномасштабной кинематикой космоса.
	
	\textbf{Мост между информацией, геометрией и космологией.}
	Механизм, описанный выше, строго выведен в Приложении Ehrenfest, где парадокс вращающегося диска разрешается путём отождествления метрического коэффициента с координационной эффективностью \(K_{\mathrm{eff}}(r)\). Та же \(K_{\mathrm{eff}}\) появляется в обобщённом законе Шеннона-Якушева (Приложение X), который устанавливает прямую эквивалентность между передачей информации, энергией и геометрией:
	\begin{equation}
		W = K_{\mathrm{eff}} \cdot I_{\text{канал}} \cdot \eta,
	\end{equation}
	где \(W\) — координированная работа (произведённый смысл), \(I_{\text{канал}}\) — переданные данные, а \(\eta\) — фактор эффективности времени. В пределе идеального координационного словаря это соотношение сводится к информационному аналогу \(E = mc^2\). Таким образом, универсальный закон масштабирования связывает теорию информации, геометрию и космологию: тот же показатель \(\beta = 2/3\) управляет потоком информации в словаре, кривизной пространства вокруг вращающегося диска и динамикой вращения Вселенной.
	
	\subsection{Простые числа как координационная лестница}
	
	Распределение простых чисел даёт независимое, чисто математическое подтверждение универсального закона ошибки. При естественной гипотезе \(K_{\mathrm{eff}}(p_n) \propto p_n\) относительное отклонение \(n\)-го простого числа от классического приближения Россера \(R_n\) должно масштабироваться как \(\varepsilon_n = |p_n - R_n|/p_n \propto p_n^{-2/3}\). Численный анализ первых \(5\times10^5\) простых чисел (Приложение~PrimeN) показывает, что локальный логарифмический наклон \(\gamma(n)\) монотонно сходится к \(2/3\) при \(n\to\infty\) (Таблица~\ref{tab:prime-slopes}).
	
	\begin{table}[H]
		\centering
		\caption{Сходимость локального показателя \(\gamma\) к \(\beta = 2/3\).}
		\label{tab:prime-slopes}
		\begin{tabular}{c c}
			\toprule
			Диапазон \(n\) & \(\gamma\) \\
			\midrule
			\(6\)--\(50\,000\)    & 0.6894 \\
			\(50\,001\)--\(100\,000\) & 0.6775 \\
			\(100\,001\)--\(150\,000\)& 0.6732 \\
			\(150\,001\)--\(200\,000\)& 0.6712 \\
			\(200\,001\)--\(250\,000\)& 0.6698 \\
			\(250\,001\)--\(300\,000\)& 0.6689 \\
			\(300\,001\)--\(350\,000\)& 0.6682 \\
			\(350\,001\)--\(400\,000\)& 0.6677 \\
			\(400\,001\)--\(450\,000\)& 0.6674 \\
			\(450\,001\)--\(500\,000\)& 0.6670 \\
			\bottomrule
		\end{tabular}
	\end{table}
	
	Из алгебраического цикла YUCT получается \textbf{беспараметрическая асимптотическая поправка} к формуле Россера:
	
	\begin{equation}
		\boxed{ p_n \approx R_n - \frac{S_{\mathrm{even}}}{2}\, n^{1-\beta}\,\ln n },
		\qquad \beta = \frac{2}{3},\quad S_{\mathrm{even}}=\frac{4}{5}.
		\label{eq:prime-yuct-theory}
	\end{equation}
	
	Эта поправка не содержит эмпирических констант. Остаточные осцилляции истинных простых чисел вокруг этого тренда следуют каскаду \textbf{фазовых переходов вакуумной решётки} на координационных глубинах
	\[
	N_f^{(k)} = 80.0 + (k-1)\cdot 16.5, \qquad k=1,2,\dots,
	\]
	где шаг \(\Delta N = 16.5\) выведен из числа фермионов Вейля \(L_0=96\) и инварианта чётного сектора \(\Delta N = L_0/6 + S_{\mathrm{even}}/2\). Знак поправки меняется на каждом пороге, как описано универсальным тригонометрическим гейтом
	\[
	\operatorname{sign\_gate}(n) = \operatorname{sign}\!\left[\sin\!\left(\frac{\pi}{16.5}\bigl(\log_q n - 80.0\bigr)\right)\right],
	\qquad q = \left(\frac{3}{2}\right)^{1/3}.
	\]
	Этот гейт — тот же фазовый оператор, который стабилизирует решётку простых чисел и непосредственно унаследован от компактного слоя \(F_{18}\) 19-мерного координационного многообразия. Его периодичность является прямым следствием алгебраических констант \(S_{\mathrm{odd}}\) и \(S_{\mathrm{even}}\). Таким образом, распределение простых чисел — это не случайный процесс, а детерминированное координационное явление, управляемое теми же законами, которые определяют массы элементарных частиц и геометрию пространства.
	
	\subsection{Связь с гипотезой математической вселенной}
	
	Гипотеза математической вселенной (MUH) Макса Тегмарка утверждает, что физическая реальность идентична математической структуре, и что все непротиворечивые математические структуры существуют физически \cite{Tegmark2014}. Хотя это элегантно решает вопрос «Почему математика?», остаётся фундаментальный пробел: \emph{почему мы наблюдаем именно эту математическую структуру, а не какую-либо другую?} Обычный ответ — антропный отбор — описателен, но не объясняет.
	
	YUCT даёт количественный ответ: \textbf{математическая структура физически реализуется тогда и только тогда, когда её элементы могут быть скоординированы с конечной эффективностью \(K_{\mathrm{eff}} \ge 1\) в соответствии с универсальным законом ошибки \eqref{eq:errorlaw}.} Распределение простых чисел удовлетворяет этому условию, поскольку подчиняется тем же законам масштабирования, что и физические системы. Напротив, чисто случайные последовательности или структуры с несовместимым масштабированием были бы нескоординированными и поэтому не могут быть реализованы как физические состояния.
	
	Конкретно, планковская граница \(N_f^{\max} = 382.0\) (Приложение~PrimeN) определяет резкий предел, за которым простые числа перестают быть различимыми с помощью любого физического словаря. Этот предел не произволен; он следует из того же алгебраического цикла, который определяет массовую лестницу и космологическую постоянную. Таким образом, YUCT превращает MUH из метафизического постулата в опровержимую теорию: любая предложенная математическая структура должна иметь конечную координационную глубину \(N_f < 382.0\) и должна демонстрировать масштабирование, совместимое с \(\beta = 2/3\), чтобы быть физически реализованной.
	
	Сходимость морфологии спиралей, статистики простых чисел и квантовой эмуляции (Приложение~PrimeN, раздел 11) на одном и том же наборе констант демонстрирует, что различие между «математическим» и «физическим» является артефактом человеческой абстракции. В реальности и то, и другое управляется единым координационным полем \(\Psi_{MN}\), действующим через протокол YPSDC. Детерминистический квантовый эмулятор из Приложения~PrimeN, построенный исключительно из матриц Штерна-Броко, даёт наглядную иллюстрацию: суперпозиция, запутанность и гейты CNOT — это обычные арифметические операции над рациональными координатами — это координированные состояния, а не таинственные неклассические явления.
	
	Таким образом, YUCT не только подтверждает дух MUH, но и даёт недостающий механизм для отбора физически реализуемых структур и устраняет необходимость в отдельной квантовой онтологии. Законы математики и законы физики — две стороны одной координационной решётки, жёсткость которой закодирована в алгебраическом цикле \(\Sodd + \Seven = 2\), \(\Sodd/\Seven = 3/2\) и \(\beta = 2/3\).
	
	\subsection{Синтез: единство информации, геометрии, математики и космологии через единый показатель}
	
	Предыдущие разделы показали, что универсальный показатель \(\beta = 2/3\) появляется в замечательном разнообразии контекстов:
	
	\begin{itemize}
		\item \textbf{Физика элементарных частиц:} массовая лестница фермионов следует \(m_f = m_e q^{N_f}\) с \(q = (3/2)^{1/3}\), непосредственно выведенным из \(\beta = 2/3\) и алгебраического цикла \(S_{\mathrm{odd}}+S_{\mathrm{even}}=2\).
		\item \textbf{Атомная и молекулярная физика:} энергии связей (HF–HI) масштабируются как \(E \propto r^{-0.682}\), что неотличимо от \(\beta = 2/3\).
		\item \textbf{Биология:} частоты ошибок репликации ДНК, метаболическое масштабирование и частоты мутаций подчиняются \(\varepsilon \propto K_{\mathrm{eff}}^{-2/3}\).
		\item \textbf{Геофизика:} \(b\)-значение Гутенберга-Рихтера удовлетворяет \(b = 1/\beta = 3/2\).
		\item \textbf{Гидродинамика и морфология спиралей:} отношение шага к радиусу любой стационарной спирали фиксировано как \(P/r = q \pi_{\mathrm{coord}} - S_{\mathrm{even}}\beta\), объединяя аммонитов, ДНК, ураганы и галактические рукава.
		\item \textbf{Теория чисел:} распределение простых чисел асимптотически сходится к базовой линии Россера с поправкой, пропорциональной \(\Seven n^{1-\beta}\ln n\), и локальный логарифмический наклон стремится к \(\beta = 2/3\).
		\item \textbf{Квантовая механика:} детерминистическая эмуляция суперпозиции, запутанности и гейтов CNOT достигается с помощью матриц Штерна-Броко с тем же периодом фазы \(\Delta N = 16.5\), выведенным из \(L_0=96\) и \(\Seven\).
		\item \textbf{Теория информации:} обобщённый закон Шеннона-Якушева \(W = K_{\mathrm{eff}} \cdot I_{\mathrm{канал}} \cdot \eta\) устанавливает прямой информационный аналог \(E = mc^2\).
		\item \textbf{Геометрия и космология:} эффективная метрика вращающегося диска имеет вид \(d\ell^2 = dr^2 + r^2 K_{\mathrm{eff}}(r)^2 d\phi^2\), а глобальный период вращения Вселенной \(T_{\mathrm{rot}} \approx 2.3\times10^{14}\) лет следует из тех же констант.
	\end{itemize}
	
	В каждом случае управляющее уравнение одно и то же:
	
	\begin{equation}
		\boxed{ \varepsilon = \kappa_c \, \alpha \, K_{\mathrm{eff}}^{-2/3} },
		\qquad \kappa_c = \frac{1}{3},\quad \beta = \frac{2}{3},
	\end{equation}
	
	где \(K_{\mathrm{eff}}\) интерпретируется в соответствии с областью: как отношение энтропии действий к энтропии индексов (информация), как отношение масс (физика частиц), как метаболическая эффективность (биология), как координация напряжений (геофизика), как координационная глубина (теория чисел) или как метрический коэффициент (геометрия). Универсальность показателя возникает потому, что все эти системы являются проекциями одного и того же координационного поля \(\Psi_{MN}\), действующего через протокол YPSDC: предварительно распределённый словарь, активируемый коротким индексом.
	
	Таким образом, YUCT раскрывает, что кажущееся разнообразие законов природы является следствием нашего ограниченного восприятия. На фундаментальном уровне существует только один закон: закон координации. Информация, энергия, геометрия и даже структура чистой математики — это разные грани единого координационного процесса. Алгебраический цикл \(\Sodd + \Seven = 2\), \(\Sodd/\Seven = 3/2\) и \(\beta = 2/3\) обеспечивает жёсткий скелет этой единой реальности, и его эмпирическое подтверждение на более чем 50 порядках величины демонстрирует, что это единство — не философская спекуляция, а измеримый факт природы.
	
	% ============================================================
	% РИСУНОК: Unified Beta
	% ============================================================
	\begin{figure}[H]
		\centering
		\begin{tikzpicture}[
			node distance=1.2cm,
			dom/.style={rectangle, draw=blue!70, fill=blue!4, thick,
				minimum width=2.8cm, minimum height=0.7cm, rounded corners, align=center, font=\small},
			center/.style={rectangle, draw=gold!80!black, fill=gold!15, thick,
				minimum width=3.2cm, minimum height=1.2cm, rounded corners, align=center, font=\bfseries\large},
			arrow/.style={->, >=Stealth, thick, color=darkblue!70},
			label/.style={font=\scriptsize\itshape, align=center}
			]
			
			% Центральный узел — универсальный закон
			\node[center] (law) at (0,0) {\(\displaystyle \varepsilon = \kappa_c \alpha K_{\mathrm{eff}}^{-2/3}\)\\[2pt]
				\small\mdseries \(\beta = 2/3,\; \kappa_c = 1/3\)};
			
			% Домены (расположены по кругу)
			\node[dom] (info) at (4.2, 2.8) {Теория информации\\Приложение X};
			\node[dom] (particle) at (5.2, 0.0) {Физика частиц\\Массовая лестница};
			\node[dom] (biology) at (4.2, -2.8) {Биология\\ДНК, метаболизм};
			\node[dom] (geo) at (0.0, -4.5) {Геофизика\\Землетрясения};
			\node[dom] (spiral) at (-4.2, -2.8) {Спиральные структуры\\14 систем};
			\node[dom] (number) at (-5.2, 0.0) {Теория чисел\\Простые числа};
			\node[dom] (quantum) at (-4.2, 2.8) {Квантовая механика\\Штерн-Броко};
			\node[dom] (cosmo) at (0.0, 4.5) {Космология\\Вращение, инфляция};
			
			% Стрелки от закона к доменам
			\draw[arrow] (law) -- (info);
			\draw[arrow] (law) -- (particle);
			\draw[arrow] (law) -- (biology);
			\draw[arrow] (law) -- (geo);
			\draw[arrow] (law) -- (spiral);
			\draw[arrow] (law) -- (number);
			\draw[arrow] (law) -- (quantum);
			\draw[arrow] (law) -- (cosmo);
			
			% Подписи к стрелкам
			\node[label, right] at (2.1, 1.6) {\(K_{\mathrm{eff}} = H(\mathcal A)/H(\mathcal K)\)};
			\node[label, right] at (3.0, 0.0) {\(K_{\mathrm{eff}} \propto m\)};
			\node[label, right] at (2.1, -1.6) {\(K_{\mathrm{eff}}\) = эффективность репарации};
			\node[label, below] at (0.0, -2.5) {\(K_{\mathrm{eff}}\) = координация напряжений};
			\node[label, left] at (-2.1, -1.6) {\(K_{\mathrm{eff}}\) = условие устойчивости};
			\node[label, left] at (-3.0, 0.0) {\(K_{\mathrm{eff}} \propto p_n\)};
			\node[label, left] at (-2.1, 1.6) {\(K_{\mathrm{eff}}\) = размер словаря};
			\node[label, above] at (0.0, 2.5) {\(K_{\mathrm{eff}}\) = метрический коэффициент};
			
			% Внешняя рамка с подписью
			\draw[gray!30, thick, rounded corners] (-6.9,-5.5) rectangle (6.9,5.8);
			\node[font=\small\bfseries, anchor=north west] at (-5.0,6.5) 
			{Универсальный закон координации во всех областях};
			
		\end{tikzpicture}
		\caption{Один и тот же показатель \(\beta = 2/3\) управляет теорией информации, физикой частиц, биологией, геофизикой, гидродинамикой, теорией чисел, квантовой механикой и космологией. В каждой области \(K_{\mathrm{eff}}\) принимает конкретную интерпретацию, но функциональная форма \(\varepsilon = \kappa_c \alpha K_{\mathrm{eff}}^{-2/3}\) остаётся инвариантной. Эта универсальность является признаком единого координационного поля \(\Psi_{MN}\).}
		\label{fig:unified-beta}
	\end{figure}
	
	\section{Проверяемые предсказания}
	\label{sec:predictions}
	
	Основываясь на законе \(\eps = \kappac\,\alpha\,\Keff^{-2/3}\) с \(\kappac = 1/3\), мы предлагаем следующие эксперименты:
	
	\begin{enumerate}
		\item \textbf{Ферментативный катализ:} для серии мутантов фермента измерить \(\Keff\) (например, по числу конформационных состояний) и каталитическую эффективность \(k_{\text{cat}}/K_M\). График \(\log(k_{\text{cat}}/K_M)\) от \(\log \Keff\) должен иметь наклон \(2/3\). Более детальные биологические предсказания даны в Приложениях D и E.
		\item \textbf{Нейронное кодирование:} в нейронной популяции оценить \(\Keff\) из взаимной информации между стимулом и ответом, делённой на энтропию ответа. Порог различения (частота ошибок) должен масштабироваться как \(\Keff^{-2/3}\).
		\item \textbf{Финансовые рынки:} используя высокочастотные данные, вычислить \(\Keff\) из ликвидности стакана заявок и спреда bid-ask. Волатильность (стандартное отклонение доходностей) должна следовать закону \(\sigma \propto \Keff^{-2/3}\).
		\item \textbf{Космическая структура:} среднеквадратичная флуктуация плотности \(\sigma_8\) должна быть связана с координационной эффективностью тёмной материи, оцененной по числу независимых областей в ранней Вселенной. Это даёт проверку согласованности с наблюдениями реликтового излучения (см. основной текст и Приложение G).
		\item \textbf{Вращающийся сверхпроводящий диск (Приложение Ehrenfest):} для диска из высокотемпературного сверхпроводника критическая угловая скорость, при которой он разрушается, должна различаться между нормальным и сверхпроводящим состояниями, поскольку координационная эффективность материала \(K_{\mathrm{eff,mat}}\) меняется. Это даёт прямую лабораторную проверку зависимости геометрии от материала.
	\end{enumerate}
	
	\section{Ограничения и оговорки}
	\label{sec:limitations}
	
	Применение универсального масштабного закона требует осторожности:
	\begin{itemize}
		\item Определение \(\Keff\) требует чёткой идентификации словаря и пространства действий; во многих системах это нетривиально.
		\item Предфактор \(\kappac = 1/3\) точен для триадной структуры. В YUCT \(\alpha\) является фиксированной константой, выведенной в Приложении Y, раздел Y.5.2; это не свободный параметр. Однако на практике, когда \(\Keff\) оценивается лишь приблизительно, эффективное значение \(\alpha\), извлекаемое из данных, может казаться переменным. Это указывает на необходимость более точных оценок \(\Keff\), а не на фундаментальную изменчивость \(\alpha\).
		\item Показатель \(2/3\) строго выведен из трёхмерности пространства (Теорема 2) и условия связности (Приложение Y), но отображение \(\Keff\) на наблюдаемые величины может требовать дополнительных предположений.
		\item Примеры в таблице~\ref{tab:data} включают системы, где ошибка не является непосредственно интересующей величиной (например, реликтовое излучение); требуется аккуратное отображение.
	\end{itemize}
	
	Тем не менее, сходимость данных из более чем 50 порядков величины убедительно свидетельствует о подлинном универсальном явлении.
	
	\section{Заключение}
	Мы представили пересмотренную, математически непротиворечивую формулировку универсального закона масштабирования ошибок в рамках YUCT. Показатель \(\beta = 2/3\) является точным и следует из трёхмерности координационного пространства (Теорема 2 основного текста) и условия связности иерархии (Приложение Y), в то время как предфактор \(\kappac = 1/3\) отражает триадную структуру. Эмпирическое значение \(0.67 \pm 0.03\) подтверждает это предсказание. Закон объединяет широкий спектр степенных законов в сложных системах и даёт проверяемые предсказания. Будущие работы должны быть сосредоточены на точных измерениях \(\Keff\) в различных областях и на расширении теоретического вывода с учётом предфактора \(\alpha\).
	
	\paragraph*{Доступность данных}
	Все компиляции данных и скрипты анализа доступны по адресу \url{https://github.com/Alexey-Yakushev-YUCT/YPSDC}.
	
	\paragraph*{История версий}
	\begin{itemize}
		\item v1.0: Первоначальная версия с противоречивыми определениями.
		\item v2.0: Исправленные определения, добавлены примеры валидации.
		\item v2.2: Расширение на атомные и биологические масштабы.
		\item v3.0: Пересмотрен вывод с использованием Теоремы 2, замена эмпирического \(0.67\) на точное \(2/3\), уточнение ограничений, добавлены ссылки на Приложения P, Y, Ferra1.2.
		\item v3.1: Дальнейшие уточнения: добавлен явный вывод множителя 2 из Приложения Y, удалён вводящий в заблуждение пример распада нейтрона, связь с моделью Изинга представлена как независимое подтверждение, добавлены ссылки на Приложения D, E, F для операционализации, исправлено обсуждение \(\alpha\) с указанием его статуса фиксированной константы.
		\item v3.2: Добавлен раздел об универсальной координации в различных областях, включая спиральные структуры и простые числа; даны ссылки на Приложения Ehrenfest и X; добавлены соответствующие библиографические ссылки.
		\item v3.3: Добавлен подраздел-синтез, явно объединяющий все области, и визуальная сводка (Рис.~\ref{fig:unified-beta}), показывающая центральную роль \(\beta = 2/3\) в теории информации, физике частиц, биологии, геофизике, гидродинамике, теории чисел, квантовой механике и космологии.
	\end{itemize}
	
	\begin{thebibliography}{99}
		\bibitem{Yakushev2026}
		A. V. Yakushev, \emph{Yakushev Unified Coordination Theory (YUCT)}, Zenodo, \url{https://doi.org/10.5281/zenodo.18444599} (2026).
		\bibitem{AppendixC1}
		A. V. Yakushev, Приложение C1: Экспериментальная проверка универсального закона масштабирования в химических связях (HF–HI), в \emph{YUCT}, Zenodo (2026).
		\bibitem{AppendixD}
		A. V. Yakushev, Приложение D: Биологические предсказания YUCT – проверяемые гипотезы, в \emph{YUCT}, Zenodo (2026).
		\bibitem{AppendixE}
		A. V. Yakushev, Приложение E: Координационная генетика – принципы YUCT в молекулярной биологии, в \emph{YUCT}, Zenodo (2026).
		\bibitem{AppendixF}
		A. V. Yakushev, Приложение F: Применение YUCT к экономике, в \emph{YUCT}, Zenodo (2026).
		\bibitem{AppendixG}
		A. V. Yakushev, Приложение G: Объединённая теория координации Якушева – фундаментальное решение проблемы квантовой гравитации, в \emph{YUCT}, Zenodo (2026).
		\bibitem{AppendixP}
		A. V. Yakushev, Приложение P: Температурная зависимость гравитации, в \emph{YUCT}, Zenodo (2026).
		\bibitem{AppendixY}
		A. V. Yakushev, Приложение Y: Математические основы YUCT – вывод из первых принципов, в \emph{YUCT}, Zenodo (2026).
		\bibitem{AppendixFerra}
		A. V. Yakushev, Приложение Ferra1.2: Универсальные координационные константы для упорядоченных и неупорядоченных фаз, в \emph{YUCT}, Zenodo (2026).
		\bibitem{AppendixPrimeN}
		A. V. Yakushev, Приложение PrimeN: Координационная лестница простых чисел YUCT, в \emph{YUCT}, Zenodo (2026).
		\bibitem{AppendixX}
		A. V. Yakushev, Приложение X: Обобщённая теория Шеннона, закон стационарной ошибки и \(K_{\mathrm{eff}}\)-зависимый предел Белла, в \emph{YUCT}, Zenodo (2026).
		\bibitem{AppendixEhrenfest}
		A. V. Yakushev, Приложение Ehrenfest: Координационная интерпретация парадокса вращающегося диска и возникновение неевклидовой геометрии, в \emph{YUCT}, Zenodo (2026).
		\bibitem{El-Showk2014}
		S. El-Showk et al., \emph{Phys. Rev. D} \textbf{90}, 105001 (2014).
		\bibitem{ammonites}
		D. K. Jacobs, Morphology and the evolution of ammonite shell coiling, \emph{Paleobiology} \textbf{16}, 283 (1990).
		\bibitem{watsoncrick}
		J. D. Watson and F. H. C. Crick, Molecular structure of nucleic acids, \emph{Nature} \textbf{171}, 737 (1953).
		\bibitem{pauling}
		L. Pauling and R. B. Corey, Configurations of polypeptide chains with favored orientations around single bonds, \emph{Proc. Natl. Acad. Sci. USA} \textbf{37}, 729 (1951).
		\bibitem{collagen}
		G. N. Ramachandran and G. Kartha, Structure of collagen, \emph{Nature} \textbf{176}, 593 (1955).
		\bibitem{phyllotaxis}
		P. H. Rivier and A. M. Turing, The chemical basis of morphogenesis, \emph{Phil. Trans. R. Soc. B} \textbf{237}, 37 (1952).
		\bibitem{cholesteric}
		P. G. de Gennes and J. Prost, \emph{The Physics of Liquid Crystals}, Oxford University Press, 2nd ed., 1993.
		\bibitem{helical-mof}
		B. Moulton and M. J. Zaworotko, From molecules to crystal engineering: supramolecular isomerism and polymorphism in coordination polymers, \emph{Chem. Rev.} \textbf{101}, 1629 (2001).
		\bibitem{ocean-eddies}
		A. R. Robinson (ed.), \emph{Eddies in Marine Science}, Springer, 1983.
		\bibitem{hurricanes}
		K. Emanuel, The physics of tropical cyclones, \emph{Annu. Rev. Fluid Mech.} \textbf{35}, 1 (2003).
		\bibitem{tornadoes}
		R. Davies-Jones, A review of supercell and tornado dynamics, \emph{Atmos. Res.} \textbf{158--159}, 274 (2015).
		\bibitem{nuclear-tracks}
		R. L. Fleischer, P. B. Price, and R. M. Walker, \emph{Nuclear Tracks in Solids: Principles and Applications}, University of California Press, 1975.
		\bibitem{binney-tremaine}
		J. Binney and S. Tremaine, \emph{Galactic Dynamics}, Princeton University Press, 2nd ed., 2008.
		\bibitem{ppdisk}
		S. M. Andrews and D. J. Wilner, The structure of protoplanetary disks, \emph{Annu. Rev. Astron. Astrophys.} \textbf{55}, 471 (2017).
		\bibitem{beck}
		R. Beck, Magnetic fields in galaxies, \emph{Space Sci. Rev.} \textbf{99}, 243 (2001).
		\bibitem{Tegmark2014}
		Max Tegmark, \emph{Our Mathematical Universe: My Quest for the Ultimate Nature of Reality}, Knopf, 2014.
	\end{thebibliography}
	
\end{document}